\documentclass[10pt,a4paper]{amsart}
\usepackage{amsmath,amssymb,amsthm,graphicx}
\usepackage{verbatim}
\newenvironment{reference}{\comment}{\endcomment}
%\newenvironment{reference}{}{}
\newenvironment{slogan}{\comment}{\endcomment}
\newenvironment{history}{\comment}{\endcomment}

% For commutative diagrams we use Xy-pic
\usepackage[all]{xy}

% We use 2cell for 2-commutative diagrams.
\xyoption{2cell}
% Theorem environments.
%
\theoremstyle{plain}
\newtheorem{theorem}[subsection]{Theorem}
\newtheorem{proposition}[subsection]{Proposition}
\newtheorem{lemma}[subsection]{Lemma}

\theoremstyle{definition}
\newtheorem{definition}[subsection]{Definition}
\newtheorem{example}[subsection]{Example}
\newtheorem{exercise}[subsection]{Exercise}
\newtheorem{situation}[subsection]{Situation}

\theoremstyle{remark}
\newtheorem{remark}[subsection]{Remark}
\newtheorem{remarks}[subsection]{Remarks}

\numberwithin{equation}{subsection}

% Macros
%
\def\lim{\mathop{\mathrm{lim}}\nolimits}
\def\colim{\mathop{\mathrm{colim}}\nolimits}
\def\Spec{\mathop{\mathrm{Spec}}}
\def\Hom{\mathop{\mathrm{Hom}}\nolimits}
\def\Ext{\mathop{\mathrm{Ext}}\nolimits}
\def\SheafHom{\mathop{\mathcal{H}\!\mathit{om}}\nolimits}
\def\SheafExt{\mathop{\mathcal{E}\!\mathit{xt}}\nolimits}
\def\Sch{\mathit{Sch}}
\def\Mor{\mathop{\mathrm{Mor}}\nolimits}
\def\Ob{\mathop{\mathrm{Ob}}\nolimits}
\def\Sh{\mathop{\mathit{Sh}}\nolimits}
\def\NL{\mathop{N\!L}\nolimits}
\def\CH{\mathop{\mathrm{CH}}\nolimits}
\def\proetale{{pro\text{-}\acute{e}tale}}
\def\etale{{\acute{e}tale}}
\def\QCoh{\mathit{QCoh}}
\def\Ker{\mathop{\mathrm{Ker}}}
\def\Im{\mathop{\mathrm{Im}}}
\def\Coker{\mathop{\mathrm{Coker}}}
\def\Coim{\mathop{\mathrm{Coim}}}

% Boxtimes
%
\DeclareMathSymbol{\boxtimes}{\mathbin}{AMSa}{"02}

%
% Macros for moduli stacks/spaces
%
\def\QCohstack{\mathcal{QC}\!\mathit{oh}}
\def\Cohstack{\mathcal{C}\!\mathit{oh}}
\def\Spacesstack{\mathcal{S}\!\mathit{paces}}
\def\Quotfunctor{\mathrm{Quot}}
\def\Hilbfunctor{\mathrm{Hilb}}
\def\Curvesstack{\mathcal{C}\!\mathit{urves}}
\def\Polarizedstack{\mathcal{P}\!\mathit{olarized}}
\def\Complexesstack{\mathcal{C}\!\mathit{omplexes}}
% \Pic is the operator that assigns to X its picard group, usage \Pic(X)
% \Picardstack_{X/B} denotes the Picard stack of X over B
% \Picardfunctor_{X/B} denotes the Picard functor of X over B
\def\Pic{\mathop{\mathrm{Pic}}\nolimits}
\def\Picardstack{\mathcal{P}\!\mathit{ic}}
\def\Picardfunctor{\mathrm{Pic}}
\def\Deformationcategory{\mathcal{D}\!\mathit{ef}}
\usepackage{cleveref}
\crefname{originallemma}{original lemma}{original lemmas}
\Crefname{originallemma}{Original lemma}{Original lemmas}
\makeatletter
\newlabel{algebra-lemma-pullback-module}{{02M1}{source}{}{sourcepointer-1}{}}
\newlabel{algebra-lemma-pullback-module@cref}{{[originallemma][0][]02M1}{[1][0][]source}{}{}{}}
\makeatother
\crefname{originalremark}{original remark}{original remarks}
\Crefname{originalremark}{Original remark}{Original remarks}
\makeatletter
\newlabel{decent-spaces-remark-functoriality-henselian-local-ring}{{0EPL}{source}{}{sourcepointer-2}{}}
\newlabel{decent-spaces-remark-functoriality-henselian-local-ring@cref}{{[originalremark][0][]0EPL}{[1][0][]source}{}{}{}}
\makeatother
\crefname{originalreference}{original reference}{original references}
\Crefname{originalreference}{Original reference}{Original references}
\makeatletter
\newlabel{lemma-Gm-torsor}{{0EQG}{source}{}{sourcepointer-3}{}}
\newlabel{lemma-Gm-torsor@cref}{{[originalreference][0][]0EQG}{[1][0][]source}{}{}{}}
\makeatother
\crefname{originalreference}{original reference}{original references}
\Crefname{originalreference}{Original reference}{Original references}
\makeatletter
\newlabel{situation-setup}{{02QL}{source}{}{sourcepointer-4}{}}
\newlabel{situation-setup@cref}{{[originalreference][0][]02QL}{[1][0][]source}{}{}{}}
\makeatother
\crefname{originalsection}{original section}{original sections}
\Crefname{originalsection}{Original section}{Original sections}
\makeatletter
\newlabel{spaces-over-fields-section-Weil-divisors}{{0ENI}{source}{}{sourcepointer-5}{}}
\newlabel{spaces-over-fields-section-Weil-divisors@cref}{{[originalsection][0][]0ENI}{[1][0][]source}{}{}{}}
\makeatother
\crefname{originalsection}{original section}{original sections}
\Crefname{originalsection}{Original section}{Original sections}
\makeatletter
\newlabel{spaces-over-fields-section-c1}{{0ENV}{source}{}{sourcepointer-6}{}}
\newlabel{spaces-over-fields-section-c1@cref}{{[originalsection][0][]0ENV}{[1][0][]source}{}{}{}}
\makeatother
\title{The divisor of a meromorphic section pulls back to the divisor of the corresponding rational function on the trivialising torsor}
\author{The Stacks Project Authors}
\date{stacks-project:spaces-chow.tex:local-snapshot-20261012; GNU FDL 1.2 or later}
\begin{document}
\maketitle
\noindent\textit{Editorial note.} All delivered mathematical text is the author's own, copied from the pinned Stacks Project chapter file spaces-chow.tex; no statement or proof step was rewritten, reordered or removed.\par\medskip
\noindent\textit{Editorial note.} The counted proof is the author's complete argument: a classification of the prime divisors of the torsor followed by a coefficient comparison through lengths over henselizations and the identification of the maximal ideals\par\medskip
\noindent\textit{Editorial note.} The label lemma-Gm-torsor-divisor-meromorphic-section names the counted statement, and the proof environment immediately following it in the source is delivered verbatim as its proof.\par\medskip
\noindent\textit{Editorial note.} The counted statement is carried by the author's own lemma environment in the chapter source.\par\medskip
\noindent\textit{Allowed standard prerequisites (editorial).}\begin{itemize}
\item Divisors of meromorphic sections of an invertible sheaf and the compatibility between the two constructions of divisors, via rational functions and via lengths
\item The trivialising torsor of an invertible sheaf on an integral good algebraic space, its flatness and the identification of the maximal ideals of the henselizations of the local rings of the torsor and of the base
\item The length criterion for the order of a section of an invertible sheaf along a prime divisor using the total ring of fractions of the henselization, and flat base change of lengths
\end{itemize}\medskip
\ifdefined\hypertarget\hypertarget{sourcepointer-1}{}\fi
\noindent\textit{Source reference.} Original lemma 02M1: sources/stacks-chow/tags.txt, source lines 3428--3428. Prerequisite pointer invoked at proof line(s) 2547; it names a standard fact of a sibling chapter that the counted argument uses, and every step of the counted proof built on it is retained verbatim in the delivered body.\par\medskip
\ifdefined\hypertarget\hypertarget{sourcepointer-2}{}\fi
\noindent\textit{Source reference.} Original remark 0EPL: sources/stacks-chow/tags.txt, source lines 18381--18381. Prerequisite pointer invoked at proof line(s) 2517; it names a standard fact of a sibling chapter that the counted argument uses, and every step of the counted proof built on it is retained verbatim in the delivered body.\par\medskip
\ifdefined\hypertarget\hypertarget{sourcepointer-3}{}\fi
\noindent\textit{Source reference.} Original reference 0EQG: sources/stacks-spaces-chow/pinned-original-source.tex, source lines 2385--2385. Prerequisite pointer invoked at proof line(s) 2482; it names a standard fact of the same chapter that the counted argument uses, and every step of the counted proof built on it is retained verbatim in the delivered body.\par\medskip
\ifdefined\hypertarget\hypertarget{sourcepointer-4}{}\fi
\noindent\textit{Source reference.} Original reference 02QL: sources/stacks-spaces-chow/pinned-original-source.tex, source lines 264--264. Prerequisite pointer invoked at proof line(s) 2478; it names a standard fact of the same chapter that the counted argument uses, and every step of the counted proof built on it is retained verbatim in the delivered body.\par\medskip
\ifdefined\hypertarget\hypertarget{sourcepointer-5}{}\fi
\noindent\textit{Source reference.} Original section 0ENI: sources/stacks-chow/tags.txt, source lines 18343--18343. Prerequisite pointer invoked at proof line(s) 2495; it names a standard fact of a sibling chapter that the counted argument uses, and every step of the counted proof built on it is retained verbatim in the delivered body.\par\medskip
\ifdefined\hypertarget\hypertarget{sourcepointer-6}{}\fi
\noindent\textit{Source reference.} Original section 0ENV: sources/stacks-chow/tags.txt, source lines 18355--18355. Prerequisite pointer invoked at proof line(s) 2496; it names a standard fact of a sibling chapter that the counted argument uses, and every step of the counted proof built on it is retained verbatim in the delivered body.\par\medskip
\begin{lemma}
\label{lemma-Gm-torsor-divisor-meromorphic-section}
In Situation \ref{situation-setup} let $X/B$ be good.
Let $\mathcal{L}$ be an invertible $\mathcal{O}_X$-module.
Assume $X$ is integral. Let $s$ be a nonzero meromorphic
section of $\mathcal{L}$. Let $q : T \to X$ be the morphism
of Lemma \ref{lemma-Gm-torsor}. Then
$$
q^*\text{div}_\mathcal{L}(s) = \text{div}_T(q^*(s))
$$
where we view the pullback $q^*(s)$
as a nonzero meromorphic function on $T$
using the isomorphism $q^*\mathcal{L} \to \mathcal{O}_T$
\end{lemma}

\begin{proof}
Observe that $\text{div}_T(q^*(s)) = \text{div}_{\mathcal{O}_T}(q^*(s))$
by the compatibility between the constructions given in
Spaces over Fields, Sections
\ref{spaces-over-fields-section-Weil-divisors} and
\ref{spaces-over-fields-section-c1}.
We will show the agreement with $\text{div}_{\mathcal{O}_T}(q^*(s))$
in this proof.
We will use all the properties of $q : T \to X$ stated in
Lemma \ref{lemma-Gm-torsor} without further mention.
Let $Z \subset T$ be a prime divisor.
Then either $Z \to X$ is dominant or $Z = q^{-1}(Z')$ for some
prime divisor $Z' \subset X$.
If $Z \to X$ is dominant, then the coefficient of $Z$ in
either side of the equality of the lemma is zero.
Thus we may assume $Z = q^{-1}(Z')$ where $Z' \subset X$
is a prime divisor.
Let $\xi' \in |Z'|$ and $\xi \in |Z|$ be the generic points.
Then we obtain a commutative diagram
$$
\xymatrix{
\Spec(\mathcal{O}_{T, \xi}^h) \ar[r]_-{c_\xi} \ar[d]_h & T \ar[d]^q \\
\Spec(\mathcal{O}_{X, \xi'}^h) \ar[r]^-{c_{\xi'}} & X
}
$$
see Decent Spaces, Remark
\ref{decent-spaces-remark-functoriality-henselian-local-ring}.
Choose a trivialization $s_{\xi'}$ of
$\mathcal{L}_{\xi'} = c_{\xi'}^*\mathcal{L}$.
Then we can use the pullback $s_\xi$ of $s_{\xi'}$ via $h$
as our trivialization of $\mathcal{L}_\xi = c_\xi^* q^*\mathcal{L}$.
Write $s/s_{\xi'} = a/b$ for $a, b \in \mathcal{O}_{X, \xi'}$ nonzerodivisors.
By definition the coefficient of $Z'$ in
$\text{div}_\mathcal{L}(s)$ is
$$
\text{length}_{\mathcal{O}_{X, \xi'}^h}(
\mathcal{O}_{X, \xi'}^h/a \mathcal{O}_{X, \xi'}^h)
-
\text{length}_{\mathcal{O}_{X, \xi'}^h}(
\mathcal{O}_{X, \xi'}^h/b \mathcal{O}_{X, \xi'}^h)
$$
Since $Z = q^{-1}(Z')$, this is also the coefficient of $Z$ in
$q^*\text{div}_\mathcal{L}(s)$. Since
$\mathcal{O}_{X, \xi'}^h \to \mathcal{O}_{T, \xi}^h$
is flat the elements $a, b$ map to nonzerodivisors in
$\mathcal{O}_{T, \xi}^h$. Thus $q^*(s)/s_\xi = a/b$ in the total quotient
ring of $\mathcal{O}_{T, \xi}^h$. By definition the coefficient of
$Z$ in $\text{div}_T(q^*(s))$ is
$$
\text{length}_{\mathcal{O}_{T, \xi}^h}(
\mathcal{O}_{T, \xi}^h/a \mathcal{O}_{T, \xi}^h)
-
\text{length}_{\mathcal{O}_{T, \xi}^h}(
\mathcal{O}_{T, \xi}^h/b \mathcal{O}_{T, \xi}^h)
$$
The proof is finished because these lengths are the same as before
by Algebra, Lemma \ref{algebra-lemma-pullback-module} and the
fact that $\mathfrak m_\xi^h = \mathfrak m_{\xi'}^h\mathcal{O}_{T, \xi}^h$
shown in Lemma \ref{lemma-Gm-torsor}.
\end{proof}

\end{document}
